\textsc{This dissertation generalises several reconstruction results}
in classical algebra to the language of monoidal categories, module categories, and monads.
We start with some brief historical remarks,
and summarise our main contributions in \cref{sec:summary} below.

Monoidal and module categories are ubiquitous in many fields of mathematics.
Briefly, a monoidal category is a category \(\cat{C}\) equipped with a tensor product functor \(\otimes\from \cat{C} \times \cat{C} \to \cat{C}\)
and a unit \(1 \in \cat{C}\), satisfying associativity and unitality conditions.
A (left) module category \(\cat{M}\) over \(\cat{C}\), in turn,
is endowed with an associative and unital action functor \(\lact \from \cat{C} \times \cat{M} \to \cat{M}\).
One motivating example of these structures comes from linear algebra.
The category \(\kVect\) of vector spaces over a field \(\Bbbk\) is monoidal,
and the right modules over a \(\Bbbk\)\hyp{}algebra \(A\) form a left \(\kVect\)\hyp{}module category.
Given \(V \in \kVect\) and a right \(A\)\hyp{}module \(M\),
the tensor product \(V \kotimes M\) is again a right \(A\)\hyp{}module,
with action given by \((v \otimes m).a \defeq v \otimes m.a\),
for all \(a \in A\), \(m \in M\), and \(v \in V\).

The areas in which the language and theory of monoidal and module categories has been applied include
quantum field theories~\cite{fuchs02:tft-rcft,day07:assoc-rcfts},
categorification in representation theory~\cite{stroppel23:categ,liu24:soerg,stroppel24:braid-soerg},
algebraic geometry~\cite{ben-zvi10:integ-drinf,ben-zvi18:integ,pascaleff18:poiss-fukay-floer},
and many aspects of the theory of Hopf algebras and tensor categories~\cite{schauenburg00:duals-r-hopf,kowalzig14:batal-vilkov,Etingof2015,fuchs22:spher-morit-serre,shimizu23:relat-serre,stroinski24:modul-tambar}.
For a survey on the role of module categories in applied category theory and computer science,
better known as actegories therein,
see~\cite{capucci22:acteg-workin-amthem}.

Just as in the category of vector spaces,
given an algebra object \(A \in \cat{C}\),
the category \(\modc{A}\) of right \(A\)\hyp{}modules in \(\cat{C}\) is naturally a left \(\cat{C}\)\hyp{}module category.
Conversely, when considering a given \(\cat{C}\)\hyp{}module category \(\cat{M}\),
it is often useful to find and study an algebra object \(A\) in \(\cat{C}\)
such that there is an equivalence \(\cat{M} \simeq \modc{A}\) of \(\cat{C}\)\hyp{}module categories.
This kind of \emph{reconstruction} process lies at the heart of \cref{sec:infinite-non-rigid-reconstruction,sec:hopf-trimodules} of the present thesis.
An early result of this kind is~\cite[Theorem~1]{ostrik03:modul-hopf},
in which a ``nice'' module category over a finite tensor category \(\cat{C}\)%
—one that is equivalent to the finite-dimensional modules over a finite-dimensional algebra—%
is expressed as an algebra object in \(\cat{C}\),
see \cref{finiterecon} of \cref{sec:infinite-non-rigid-reconstruction}.
Many generalisations and variants have appeared since;
for instance
\cite[Theorem~7.10.1]{Etingof2015},
\cite[Theorem~2.24]{douglas19},
\cite[Theorem~4.7]{mackaay19:simpl},
and \cite[Theorem~4.6]{ben-zvi18:integ}.

All of the reconstruction results cited above involve rigid monoidal categories:
those in which every object has a dual.
For example, to every finite-dimensional vector space \(V\),
there exist canonical evaluation and coevaluation maps
\(\ev \from \rd*{V} \kotimes V \to \Bbbk\) and \(\coev\from \Bbbk \to V \kotimes \rd*{V}\),
where \(\rd*{V} \defeq \mathrm{Hom}_{\Bbbk}(V, \Bbbk)\) is the dual space.
However, in~\cite[Example~2.20]{douglas19} it is shown that, in the absence of rigidity,
there exist \(\cat{C}\)\hyp{}module categories from which it is impossible to reconstruct an algebra object in \(\cat{C}\).
This forces us to consider a further generalisation—monads.

\froufrou%

\textsc{Historically, the study of} (co)monads arose out of (co)homology and homotopy theory,\marginnote{%
  The term ``monad'' first appears in print in Section~5.4 of~\cite{benabou67:introd}.
  Other names include ``triple'', ``standard con\-struc\-tion'', or ``triad''.
}
where one can use comonads to construct simplicial resolutions, see~\cite{godement58:topol,huber61:homot}.
Already noted in~\cite{huber61:homot} is the deep connection of monads with the theory of adjunctions:
every adjunction gives rise to a monad
and, in turn, every monad arises from an adjunction.
In fact, there exists a category of all adjunctions that realise a given monad.
The initial and terminal object therein%
—the Kleisli and Eilenberg–Moore category of the monad,~\cite{eilenberg65:adjoin,kleisli65:every}—%
play a pivotal role in this work.

Further applications are manifold and include
category theory~\cite{Beck1969,Street1972,blackwell89:two,lack02,arkor24},
categorical representation theory~\cite{Moerdijk2002,Bruguieres2011,stroinski24:modul-tambar,batista22:como-part},
Hopf algebras and quantum groups~\cite{Bruguieres2007,Aguiar2012,kowalzig15:cyclic,hasegawa-lemay2018:LinearHopf},
universal algebra~\cite{linton66:aspects,walters70:cat-ualg,hyland07},
formal semantics~\cite{moggi89:comput,moggi91:notion,turi97-math-op},
functional programming~\cite{wadler92:compr,wadler93:mon-fun-prog,peytonjones93:imp-fun-prog}, and
algebraic effects~\cite{benton02:monad,plotkin02:notion,plotkin09:handl,bauer15:prog-alg}.

For this thesis,
the most important aspect of monads lies in
monadic reconstruction theory%
—relating additional structure on a monad \(T\) on \(\cat{C}\)
to structure on its Eilenberg–Moore category \(\cat{C}^T\).
The connection to the theory of adjunctions
yields a canonical forgetful functor \(U^T\from \cat{C}^T \to \cat{C}\).
This relationship can thus be seen as
a generalisation of \emph{Tannaka--Krein duality}~\cite{tannaka38:ueber-dualit-grupp,krein49:princ-ual-block-alg},
in which one reconstructs a compact group from its category of representations.
Generalisations of this type of result have been proposed in a variety of different contexts,
see~\cite{%
  saavedra72:categ,%
  woronowicz88:tannak-krein,%
  deligne90:categ,%
  lurie04:tannak-dualit-geomet-stack,%
  szlachanyi09:fiber-tannak,%
  schaeppi13:tannak%
}.
A crucial feature of all of these statements is the involvement of a \emph{fibre functor}:
a forgetful strict monoidal functor to a nice underlying category,
like the category of vector spaces or bimodules over a ring.

In the monadic case,
Tannaka--Krein-type reconstruction results were obtained for bimonads~\cite{Moerdijk2002,McCrudden2002},
Hopf monads~\cite{Bruguieres2007},
as well as linearly distributive and \(*\)\hyp{}autonomous monads~\cite{pastro09:closed,pastro12:note}.
We enrich this already broad landscape of results
with new insights from the module category\hyp{}theoretic perspective,
offering a new reconstruction result for comodule monads over bimonads
in the spirit of the algebraic version of considering coideal subalgebras over bialgebras.

\section{Summary}\label{sec:summary}

\textsc{We shall now provide a short outline};
a more detailed description and introduction will be provided at the start of each respective chapter.
All original contributions of this thesis are already contained in either of the following articles or preprints.%
\marginnote{%
  The order in which the articles are listed here is that in which they were uploaded to the
  \href{https://arxiv.org/}{arXiv}.%
}

\begin{itemize}
  \item\cite{halbig24:pivot-hopf}:
  Sebastian Halbig and Tony Zorman.\ %
  Pivotality, twisted centres, and the anti\hyp{}double of a Hopf monad.\ %
  In: Theory Appl.~Categ.~41 (2024), pp.~86–149.\ issn:~1201-561\textsc{x}.
  \item\cite{halbig23:diagr-comod-monad}:
  Sebastian Halbig and Tony Zorman. Diagrammatics for Comodule Monads.\ %
  In: Appl.~Categ.~Struct.~32 (2024).\ Id/No 27, p.~17.\ %
  issn:~0927-2852.\ %
  doi:~10.1017/cbo9781139542333.
  \item\cite{halbig23:dualit-monoid-categ}:
  Sebastian Halbig and Tony Zorman.\ %
  Duality in Monoidal Categories; 2023.\ %
  arXiv:~\href{https://arxiv.org/abs/2301.03545}{2301.03545}.
  \item\cite{stroinski2024:reconstr}:
  Mateusz Stroiński and Tony Zorman.\ %
  Reconstruction of module categories in the infinite and non-rigid settings; 2024.\ %
  arXiv:~\href{https://arxiv.org/abs/2409.00793}{2409.00793}.
  \item\cite{zorman25:duoid-r-matric}:
  Tony Zorman.\ %
  Duoidal R-Matrices; 2025.\ %
  arXiv:~\href{https://arxiv.org/abs/2503.03445}{2503.03445}.
\end{itemize}
To keep a common thematic focal point, the article~\cite{coulembier25:simpl}%
—although created during the author's PhD studies—%
will not be part of this work.

\bigskip
In the interest of brevity, we shall refrain from citing either of the articles%
~\cite{%
  halbig23:dualit-monoid-categ,%
  halbig23:diagr-comod-monad,%
  halbig24:pivot-hopf,%
  stroinski2024:reconstr,%
  zorman25:duoid-r-matric%
}
beyond this introduction.

\subsection*{\Cref{chap:preliminaries}: Preliminaries}

Being an amalgam of the preliminaries of all articles,
this chapter fixes notation
and recalls basic properties of, for example,
2-categories,
monoidal and module categories,
the Drinfeld centre,
linear and abelian categories,
coends,
as well as monads.
We furthermore introduce one of the main tools for computation employed throughout—the graphical calculus.

\subsection*{\Cref{sec:trep-categories}: Duality theory for monoidal categories}

Being mainly based on~\cite{halbig23:dualit-monoid-categ},
in this chapter we study various dualities for monoidal categories and how they interact:
rigidity, tensor representability, Grothendieck–Verdier duality, and closedness.
These notions form a hierarchy,
with closedness—the least restrictive type of duality—on one end,
and rigidity on the other.
However, it is not obvious whether all inclusions are strict.
Taken together,
% Putting everything into one \cref here produces an extra ", " :(
\cref{prop:right-tensor-rep->right-gv,prop:gv->closed,prop:trep-criteria};
\cref{ex:r-cat-not-always-trep,rmk:closed-but-not-gv}; and
\cref{thm:props-of-c}
yield
\[
  \text{Rigid}
  \subsetneq \text{Tensor representable}
  \subsetneq \text{\GV}
  \subsetneq \text{Closed},
\]
answering a question of Heunen.

We then more closely investigate duality structures
on finite\hyp{}dimensional functor categories
endowed with Day convolution as a tensor product.
Intuitively, this generalises the tensor product of modules over a commutative algebra,
see \cref{ex:day-conv-as-tensor-product}.
Given a nice base category,
closedness and \GV{} duality
lift to this setting.

\begin{reptheorem}{prop:fun-cat-GV}
  Let \(\cat{C}\) be a \(\Bbbk\)\hyp{}linear hom-finite Grothendieck–Verdier category.
  Under the assumptions of \cref{hypothesis:tens-repr-mack},
  the finite-dimensional functor category \([\cat{C}, \kvect]\) is a Grothendieck–Verdier category.
\end{reptheorem}

Being motivated by representation theoretic questions,
the Cauchy completion of a category
also plays an important role for us.
For example,
the full subcategory of projective modules over a ring
is the Cauchy completion of
the full subcategory of all free modules.
We find that the Cauchy completion completely mirrors the duality behaviour of its base category.

\begin{reptheorem}{cor:pres-ref-dual-struct}
  Let \(\cat{C}^{\op}\) be a \(\Bbbk\)\hyp{}linear right closed monoidal category.
  Then \(\cat{C}^{\op}\) is right rigid (tensor representable)
  if and only if
  its Cauchy completion \(\overline{\cat{C}^{\op}}\) is right rigid (tensor representable).
  Further, \((\cat{C}^{\op}, d)\) is a right \GV{} category
  if and only if
  \((\overline{\cat{C}^{\op}}, \yo_{d})\) is.
\end{reptheorem}

\subsection*{\Cref{sec:twisted-centres}: Twisted centres}

In this chapter, mainly based on~\cite{halbig24:pivot-hopf},
we hoist \cref{thm: equiv_in_Hopf_setting}
from the Hopf-theoretic world
into that of monoidal categories.
It lays the groundwork for later monadic generalisations.

\begin{reptheorem}{thm: equivalence_in_the_rigid_setting}
  Let \(\cat{C}\) be a rigid category.
  There is a bijection between
  \begin{thmlist}[noitemsep]
    \item equivalence classes of quasi-pivotal structures on \(\cat{C}\),
    \item the Picard heap of the anti\hyp{}centre, and
    \item isomorphism classes of equivalences of module categories between the centre and the anti\hyp{}centre.
  \end{thmlist}
\end{reptheorem}

We then more closely investigate the heap structure of the anti-centre,
answering a question of Shimizu about the relationship to quasi-pivotal structures on the underlying category.

\begin{reptheorem}{thm: non_induced_pivotal_structure}
  There exists a category \(\cat{C}\)
  on which there exists a pivotal structure
  that is not induced by any element of the Picard heap of the anti\hyp{}centre of \(\cat{C}\).
\end{reptheorem}

\subsection*{\cref{sec:monadic-tannaka-reconstruction}: Monadic Tannaka--Krein reconstruction}

In this chapter,
which is mainly based on~\cite{halbig23:diagr-comod-monad,halbig24:pivot-hopf,stroinski2024:reconstr},
we study comodule monads in the sense of~\cite{Aguiar2012}.
These structures can be seen as generalising comodule algebras over a bialgebra, see \cref{ex:comodule-algebra}.
We then prove our main Tannaka--Krein reconstruction result for comodule monads
in the spirit of~\cite[Theorem~7.1]{Moerdijk2002} and~\cite[Corollary~3.13]{McCrudden2002}.

\begin{reptheorem}{thm:comodule-monad-reconstruction}
  Let \(B\) be a bimonad on the monoidal category \(\cat{C}\),
  and \(T\) a monad on a right \(\cat{C}\)\hyp{}module category \(\cat{M}\).
  Coactions of \(B\) on \(T\) are
  in bijection with
  right actions of \(\cat{C}^B\) on \(\cat{M}^T\) such that \(U^{T}\) is a strict comodule functor over \(U^{B}\).
\end{reptheorem}

In particular, we prove statements akin to Kelly's doctrinal adjunctions result
in the case of comodule and \(\cat{C}\)\hyp{}module monads%
—the latter play an important role in \cref{sec:infinite-non-rigid-reconstruction,sec:hopf-trimodules}.

\begin{reptheorem}{thm:comodule-adjunctions-correspond-to-strong-comodule-functors,porism:doctrinal-module-adjunctions}
  Let \(\stdadj\) be an oplax monoidal adjunction.
  Lifts of an ordinary adjunction \(\adj{G}{V}{\cat{M}}{\cat{N}}\) to a comodule adjunction are
  in bijection with lifts of \(V \from \cat{N} \to \cat{M}\) to a strong comodule functor.

  An adjunction \(\adj{F}{U}{\cat{M}}{\cat{N}}\) between \(\cat{C}\)\hyp{}module categories
  yields a bijection of oplax \(\cat{C}\)\hyp{}module structures on \(F\)
  and lax \(\cat{C}\)\hyp{}module structures on \(U\).
\end{reptheorem}

\subsection*{\Cref{sec:monad-twist-centr}: Monadic twisted centres}

This chapter,
mainly based on~\cite{halbig24:pivot-hopf},
provides a monadic interpretation of \cref{sec:twisted-centres}.
Based on the Drinfeld double of a Hopf monad, we develop the notion of an anti\hyp{}double,
which realises the anti\hyp{}centre of \emph{ibid} as its Eilenberg–Moore category,
and present a monadic variant of \cref{thm: equivalence_in_the_rigid_setting}.
\begin{reptheorem}{thm: equivalence_in_the_monad_setting}
  Let \(\cat{C}\) be a rigid monoidal category.
  For a Hopf monad \(H\) on \(\cat{C}\) that admits a double and anti\hyp{}double,
  the following statements are equivalent:
  \begin{thmlist}[noitemsep]
    \item the monoidal unit of \(\cat{C}\) lifts to a module over the anti-double,
    \item the double and anti\hyp{}double of \(H\) are isomorphic as comodule monads, and
    \item the double and anti\hyp{}double of \(H\) are isomorphic as monads.
  \end{thmlist}
  If \(\cat{C}\) is pivotal,
  the above statements hold if and only if \(H\) admits a pair in involution.
\end{reptheorem}

Investigating the interplay between the double and anti\hyp{}double,
we obtain a new criterion for when a rigid monoidal category is pivotal.

\begin{reptheorem}{cor: pivotal_from_central_anti_central}
  Let \(\cat{C}\) be a rigid monoidal category.
  If \(\cat{C}\) admits a central Hopf monad \(\mathfrak{D}(\cat{C})\)
  and an anti\hyp{}central comodule monad \(\mathfrak{Q}(\cat{C})\),
  then it is pivotal if and only if \(\mathfrak{D}(\cat{C}) \cong \mathfrak{Q}(\cat{C})\) as monads.
\end{reptheorem}

\subsection*{\Cref{sec:duoidal-reconstruction}: Duoidal R-matrices}

Being mainly based on~\cite{zorman25:duoid-r-matric},
this chapter introduces R\hyp{}matrices in duoidal categories,
generalising the well\hyp{}known R\hyp{}matrices for bialgebras,
and those for bimonads~\cite{Bruguieres2007}.
As it turns out, R-matrices on a suitable monad correspond to duoidal structures on its Eilenberg–Moore category.

\begin{reptheorem}{thm:r-matrices-iff-duoidal-structure}
  Let \(\cat{D}\) be a category with monoidal structures \(\circ\) and \(\bullet\),
  and \(T\) a monad on \(\cat{D}\) that has a \(\circ\)\hyp{}oplax monoidal and a \(\bullet\)\hyp{}oplax monoidal structure.
  Then quasi\-triangular structures on \(T\) are in bijection with duoidal structures on \(\cat{D}^T\).
\end{reptheorem}

We then explore the connection between normal duoidal categories
and (non-planar) linearly distributive categories~\cite{cockett97:weakl,garner16:commut}.
The latter can be seen as an analogue of Grothendieck–Verdier categories without an explicit notion of dual.
In \cref{prop:cocommutative-double opmonoidal monads-are-pastro-monads},
we relate a cocommutative version of duoidal monads
to the linearly distributive monads of~\cite{pastro12:note}.

\subsection*{\Cref{sec:infinite-non-rigid-reconstruction}: Infinite and non-rigid reconstruction}

This chapter,
mainly based on~\cite{stroinski2024:reconstr},
studies reconstruction without a fibre functor,
as done for example in~\cite[Theorem~1]{ostrik03:modul-hopf}.

Since every object \(m\) in a module category \(\cat{M}\) over \(\cat{C}\)
gives rise to a strong \(\cat{C}\)\hyp{}module functor \(\blank \lact m \from \cat{C} \to \cat{M}\),
and in good cases this functor has a right adjoint,
the resulting monad is naturally a lax \(\cat{C}\)\hyp{}module monad.
However, it is a priori not clear whether the Eilenberg–Moore category
of a (right exact) lax \(\cat{C}\)\hyp{}module monad even has a canonical \(\cat{C}\)\hyp{}module structure.
Using multicategorical techniques, we establish when such an \emph{extension} exists.

\begin{reptheorem}{extendingalways,strongembeddingalways}
  Let \(\cat{C}\) be an abelian monoidal and \(\cat{M}\) an abelian \(\cat{C}\)\hyp{}module category.
  Suppose that the action functor \(\lact \from \cat{C} \kotimes \cat{M} \to \cat{M}\) is right exact in both variables,
  and let \(T\from \cat{M}\to \cat{M}\) be a right exact lax \(\cat{C}\)\hyp{}module monad.
  Then there exists an essentially unique \(\cat{C}\)\hyp{}module structure on \(\cat{M}^{T}\),
  such that the canonical inclusion \(\iota\from \cat{M}_T \to \cat{M}^T\) is a strong \(\cat{C}\)\hyp{}module functor.
\end{reptheorem}

Using the constructed \(\cat{C}\)\hyp{}module structure,
we then establish a reconstruction and classification result
under mild additional assumptions.

\begin{reptheorem}{mainnonsenseprojective,mainnonsenseinjective,projcorrespondence}
  Assume that \(\cat{C}\) and \(\cat{M}\) have enough projectives (injectives)
  and that there exists an object \(\ell \in \cat{M}\) such that:
  \begin{itemize}[noitemsep]
    \item there is a right adjoint \(\hom{\ell,\blank}\) (left adjoint \(\cohom{\ell,\blank}\)) to \(\blank \triangleright \ell\);
    \item for \(x \in \cat{C}\) projective (injective), the object \(x \triangleright \ell\) is projective (injective);
    \item any projective (injective) object of \(\cat{M}\) is a direct summand of an object of the form \(x \triangleright \ell\), for \(x\) projective (injective).
  \end{itemize}
  Let \(T\) be the monad \(\hom{\ell, \blank\,\triangleright\,\ell}\).
  Then there is an equivalence \(\cat{M} \simeq \cat{C}^T\) of \(\cat{C}\)\hyp{}module categories,
  where the \(\cat{C}\)\hyp{}module structure of the category of \(T\)\hyp{}modules is extended from the Kleisli category.
  Furthermore, this gives rise to a bijection
  \begin{align*}
    \faktor{\{\,(\cat{M},\ell) \text{ as above}\,\}}{(\cat{M} \simeq \cat{N})}
    &\;\leftrightarrow\,
      \Bigg\{
      \begin{aligned}
        &\text{Right exact lax \(\cat{C}\)\hyp{}module}\\
        &\text{monads on \(\cat{C}\)}
      \end{aligned}
      \Bigg\}
      /\, (\cat{C}^T \simeq \cat{C}^S) \\
    (\cat{M},\ell) &\longmapsto \hom{\ell,-\lact \ell} \\
    (\cat{C}^T, T1) &\longmapsfrom T
  \end{align*}
\end{reptheorem}

\subsection*{\Cref{sec:hopf-trimodules}: Hopf trimodules}

The final chapter,
mainly based on~\cite{stroinski2024:reconstr},
further develops the non-rigid reconstruction results of \cref{sec:infinite-non-rigid-reconstruction}
in the case of several examples.

For a Hopf algebra \(B\),
the \emph{Fundamental Theorem of Hopf Modules}
says that the category \(\Tetramod[B][B]\) of Hopf modules is
naturally equivalent to
the category of vector spaces~\cite{larson69:hopf}.
Likewise, the category of \emph{Hopf trimodules} \(\Trimod\) is equivalent to that of right comodules.
In comparison to plain Hopf modules, one can generalise Hopf trimodules to the quasi-Hopf algebraic world.

If \(B\) is just a bialgebra%
—\ie, does not have an antipode—%
then we can still extract interesting structure from such modules.

\begin{reptheorem}{thm:hopf-trimodules-are-lax-monoidal-functors}
  For a bialgebra \(B\) and \(\cat{V} \defeq \Tetramod[B]\),
  there is a monoidal equivalence
  \[
    \Trimod \simeq \mathsf{LexfLax}\cat{V}\mathsf{Mod}(\cat{V}, \cat{V})
  \]
  between the category of Hopf trimodules, and
  the category of left exact finitary lax \(\cat{V}\)\hyp{}module endofunctors on \(\cat{V}\).
\end{reptheorem}

This allows us to give a categorical interpretation of
the aforementioned Fundamental Theorem of Hopf Modules.

\begin{reptheorem}{LaxIsStrongThenLeftRigid,cor:twisted-anti-iff-bestoftimes}
  The functor \(\Tetramod[B] \to \Tetramod[B][B][][B]\) corresponds to
  the inclusion of strong \(\Tetramod[B]\)\hyp{}module endofunctors into the category of lax \(\Tetramod[B]\)\hyp{}module endofunctors.
  It is an equivalence if and only if
  \(\tetramodfd[B]\) is left rigid,
  which is the case if and only if \(B\) has a twisted antipode.
\end{reptheorem}

We also give applications in the setting of semigroup algebras, see \cref{coreps,fiberfunctorthings}.
This can be seen as an example of a non-rigid category
where the reconstruction procedure described in~\cite[Chapter~7]{Etingof2015} fails.

Lastly, the fusion operators of a bimonad may be expressed in terms of its canonical module structure.

\begin{reptheorem}{delusion,diffusion}
  Let \(\adj{F}{U}{\cat{C}}{\cat{D}}\) be an oplax monoidal adjunction.
  The strong monoidal structure of \(U\) turns \(\cat{C}\) into a \(\cat{D}\)\hyp{}module category, by defining \(\blank \lact \bblank \defeq U(\blank) \otimes \bblank\).
  With respect to this \(\cat{D}\)\hyp{}module structure,
  the bimonad \(T \defeq UF\) on \(\cat{C}\) becomes an oplax \(\cat{D}\)\hyp{}module monad.
  The right fusion operator is given by
  evaluating the \(\cat{D}\)\hyp{}module coherence morphism at a free \(T\)\hyp{}module,
  and furthermore \(T\) is right Hopf
  if and only if
  it is a strong \(\cat{C}^T\)\hyp{}module monad.
\end{reptheorem}

%%% Local Variables:
%%% mode: LaTeX
%%% TeX-master: "../main"
%%% TeX-command-extra-options: "-shell-escape -fmt=prec.fmt -file-line-error -synctex=1"
%%% End:
